% Drawing 136/19-1 -- folder 136, batch 19, archive page 374 (PDF page 375).
% Transcription: à droite, un petit diagramme du plan $(i,j)$ : la ligne $j=n$ porte des $0$, et la région au-dessus est hachurée.
% Kind: small diagram of the (i, j) plane with hatched region. Method 2 (schematic redraw from the facsimile). Status: ok.
% Notes: Axes (horizontal with an arrow), the line j = n (labelled n on the left), oblique hatching strokes above it, and two 0s just below the line.
% Coordinates: 1 unit = 1 mm = 1/100 of the crop width (bbox in page fractions: [0.655, 0.685, 0.8, 0.76]).
% Redrawn by Claude (claude-opus-5-5), issue #73, 2026-11-07; not checked by a human; not compiled (no TeX in the session).
% Rights: reproduces part of the work -- behind RENDER_TRANSCRIPTS, never published by itself.
\documentclass[tikz,border=3pt]{standalone}
\usepackage{amsmath,amssymb}
\usetikzlibrary{patterns,arrows.meta}
\tikzset{pencil/.style={black!45}, struck/.style={black!35}, mark/.style={circle,draw,inner sep=1.3pt}, pt/.style={circle,fill,inner sep=0.9pt}, opt/.style={circle,draw,fill=white,inner sep=1.1pt}}
\begin{document}
\begin{tikzpicture}[x=1mm,y=1mm,line cap=round,line join=round,font=\scriptsize]
\draw (25,-60) -- (24,-10);
\draw[-{Stealth[length=2mm]}] (16,-48) -- (84,-46);
\draw (25,-33) -- (80,-33);
\node[inner sep=1pt] at (15,-36) {\tiny $n$};
\draw (30,-15) -- (38,-28);
\draw (40,-9) -- (48,-26);
\draw (55,-12) -- (62,-22);
\draw (68,-13) -- (70,-18);
\node[inner sep=1pt] at (39,-41) {\scriptsize $0$};
\node[inner sep=1pt] at (60,-40) {\scriptsize $0$};
\end{tikzpicture}
\end{document}
